\section{Consequences of the operator theorem}
\label{cons:section}

We now prove Corollary~\ref{cor:critical-laws}. The two parts of the
connection-operator conclusion play related but distinct roles.
Boundedness at criticality makes the triangle diagram finite, bringing
the standard mean-field theory into play.
The strict interval of operator boundedness above $\pc$ gives
exponential two-point decay there and supplies the stronger input
used for the extrinsic-radius law.

\begin{proof}[Proof of Corollary~\ref{cor:critical-laws}]
For every graph in the corollary, \cref{thm:main} makes
$T_{\pc}$ a bounded operator on counting-measure $\ell^2(V)$.
Tonelli's theorem identifies its nonnegative matrix product with the
triangle sum. If $\delta_v$ is the unit vector at $v$, then
\[
 \nabla_{\pc}(v)
 =\ip{\delta_v}{T_{\pc}^{3}\delta_v}
 \le\norm{T_{\pc}^{3}}_{2\to2}
 \le\norm{T_{\pc}}_{2\to2}^{3}.
\]
This proves \eqref{eq:critical-triangle}. The same implication is
recorded in \cite[Section~1]{Hutchcroft2020L2}.

\paragraph{Critical laws on simple graphs.}
First suppose that $G$ is simple. The susceptibility,
percolation-probability, and cluster-size estimates are the classical
mean-field consequences of the triangle condition
\cite{AizenmanNewman1984,BarskyAizenman1991}; their quasi-transitive
form is recorded in
\cite[Section~1, Equations~(1.1)--(1.3)]{Hutchcroft2020L2}.
For explicit applications with the precise quasi-transitive hypotheses,
we use the following later results. Under the stronger condition
$\pc<\ptwo$ available here, we may apply
\cite[Theorem~1.1 and Corollary~1.4]{Hutchcroft2022FiniteClusters}.
The corollary with $k=1$ and $p\uparrow\pc$ gives
\eqref{eq:meanfield-chi}, since subcritical clusters are finite.
The theorem at $p=\pc$ gives \eqref{eq:meanfield-size}, since
Lemma~\ref{pre:critical-finiteness} excludes infinite critical clusters.

For percolation probability, Lemma~2.1 of
\cite{Hutchcroft2022FiniteClusters}, whose hypothesis is
the triangle condition, bounds the intrinsic-radius tail uniformly
over vertices. Letting the radius tend to infinity gives
$\theta_v(p)\le C(p-\pc)$ for $p\ge\pc$: local finiteness makes an
infinite cluster equivalent to an unbounded intrinsic radius.
For the lower bound, choose one representative from each vertex
orbit, forming a finite set $F$. Theorem~2 and Propositions~11
and~23 of \cite{AntunovicVeselic2008}, with
$\beta=-\log(1-p)$, give
\[
 \sum_{u\in F}\theta_u(p)\ge c(p-\pc)
\]
for $p>\pc$ sufficiently close to $\pc$. Indeed, their bound is
linear in $\beta-\beta_c$, and this is comparable to $p-\pc$ near
$0<\pc<1$. Fix $v$ and choose a path of length $l_u$ from $v$ to
each $u\in F$. Harris association applied to the event that this
path is open and the event $\{|C_u|=\infty\}$ gives
$\theta_v(p)\ge p^{l_u}\theta_u(p)$. Hence
\[
 \theta_v(p)\ge
 \frac{\min_{u\in F}\pc^{l_u}}{|F|}
 \sum_{u\in F}\theta_u(p)
 \ge c_v(p-\pc),
\]
which proves \eqref{eq:meanfield-theta}.

For the radius estimates, it is important to distinguish the two
metrics. The intrinsic-radius upper bound under the triangle condition
originates in the lattice setting in
\cite[Theorem~1.2(ii)]{KozmaNachmias2009}; its
quasi-transitive form and the matching general lower bound are
described in \cite[Section~3]{Hutchcroft2020L2}. The extrinsic-radius
law uses the stronger condition $\pc<\ptwo$
\cite[Theorem~3.1]{Hutchcroft2020L2}. Both estimates in the exact
quasi-transitive form needed here are included in
\cite[Theorem~1.2]{Hutchcroft2022FiniteClusters}, whose hypothesis is
$\pc<\ptwo$. Set $p=\pc$ in that theorem. Its events require the
cluster radius to be finite; Lemma~\ref{pre:critical-finiteness} shows that
all critical clusters are finite almost surely, so this restriction
may be removed. The resulting estimates are
\eqref{eq:meanfield-intrinsic} and \eqref{eq:meanfield-extrinsic}.

\paragraph{Exponential connectivity.}
The definition of $\pexp$ above is equivalent to that in
\cite[Section~2]{Hutchcroft2020L2}: a positive exponential rate is
equivalent to an estimate with some positive exponent and finite
prefactor, uniformly over vertices. Theorem~2.2 of that paper gives
$\ptwo\le\pexp$. More precisely, its Remark~2.3 gives, for
$0<p<\ptwo$,
\[
 T_p(x,y)\le
 2\norm{T_p}_{2\to2}
 \exp\left(-\frac{\dist_G(x,y)}
                 {e\norm{T_p}_{2\to2}}\right)
 \qquad(x,y\in V).
\]
At $p=0$ the off-diagonal connection probabilities vanish. Together
with $\pc<\ptwo$ from \cref{thm:main}, this proves
\eqref{eq:exponential-threshold} and
\eqref{eq:exponential-connectivity} for simple graphs.

\paragraph{Subdivision and transfer of the laws.}
It remains to pass from simple graphs to the full graph class of the
corollary. Delete the loops, which change none of the observables,
and let $H=\widetilde G$ be the simple subdivided graph in
Lemma~\ref{pre:simple-reduction}. Denote the set of nonloop edge copies of
$G$ by $E_0$, so $V(H)=V(G)\sqcup E_0$. Let $d_0$ bound the number
of incident nonloop edges at a vertex of $G$, counting parallel copies.
In the coupling of that lemma, parameter $t$ on $H$ induces parameter
$p=t^2$ on $G$, old-vertex connection probabilities agree, and
$\pc(H)^2=\pc(G)$.

For an old root $v$, write $C_v^G$ and $C_v^H$ for the two coupled
clusters, and use the same superscripts for the connection and
percolation observables. Their old vertices agree. Every new vertex
in $C_v^H$ is joined by an open edge to an old neighbor in the same
cluster, and every old vertex has at most $d_0$ new neighbors.
Hence, pathwise,
\begin{equation}\label{eq:subdivision-size}
 |C_v^G|\le |C_v^H|\le(1+d_0)|C_v^G|,
 \qquad
 \theta_v^G(t^2)=\theta_v^H(t).
\end{equation}
For two old vertices, both the ambient distance and the distance along
open paths in $H$ are twice their counterparts in $G$. A new vertex
in the root cluster is at distance at most one from an old vertex in
that cluster. Consequently, for either radius,
\begin{equation}\label{eq:subdivision-radius}
 2\operatorname{rad}_{*}(C_v^G)
 \le \operatorname{rad}_{*}(C_v^H)
 \le2\operatorname{rad}_{*}(C_v^G)+1,
 \qquad *\in\{\mathrm{int},\mathrm{ext}\}.
\end{equation}
The statements include infinite values. Applied at
$t=\pc(H)=\sqrt{\pc(G)}$, these comparisons transfer the three
critical tails from $H$ to $G$. They also transfer susceptibility
and percolation probability on the appropriate sides of criticality:
if $t_c=\sqrt{\pc(G)}>0$, then
$p-\pc(G)=(t-t_c)(t+t_c)$, whose second factor stays bounded above
and away from zero near $t_c$. Multiplicative constants may change,
as allowed in the corollary.

\paragraph{Transfer of the operator threshold.}
To transfer the full range in \eqref{eq:exponential-connectivity},
we also compare the two operator thresholds. Define the incidence
operator $B:\ell^2(E_0)\to\ell^2(V(G))$ by
$B(x,e)=\one_{\{x\text{ is an endpoint of }e\}}$.
Its row sums are at most $d_0$ and its column sums are two.
Cauchy--Schwarz at each old vertex gives
$\norm{Bf}_2^2\le2d_0\norm f_2^2$, so
$\norm B_{2\to2}\le\sqrt{2d_0}$. The operator
\[
 Q:\ell^2(V(G))\longrightarrow
       \ell^2(V(G))\oplus\ell^2(E_0),\qquad
 Qf=(f,B^*f),
\]
satisfies $\norm Q_{2\to2}^{2}\le1+2d_0$.
For an old-to-new connection in $H$, a simple open path must enter the new
vertex through one of its two old endpoints. For a connection between
distinct new vertices, choose an old endpoint at each end of a simple open
path. Between these old endpoints the path induces a connection in
$G$. The union bound, together with equality for old-to-old
connections, therefore gives the entrywise inequality
\[
 T_t^H\le Q T_{t^2}^G Q^*.
\]
The product on the right is interpreted entrywise here; each entry
is a sum over at most four pairs of old endpoints.
This inequality also covers a new diagonal entry: the corresponding
diagonal entry of $B^*T_{t^2}^G B$ is at least two.
If $T_{t^2}^G$ is bounded, entrywise domination of nonnegative
kernels makes $T_t^H$ bounded with norm at most
$\norm Q^2\norm{T_{t^2}^G}$. Conversely, if $T_t^H$ is bounded,
compression to the old vertices makes $T_{t^2}^G$ bounded.
Thus boundedness is equivalent and, in that case,
\[
 \norm{T_{t^2}^G}_{2\to2}
 \le\norm{T_t^H}_{2\to2}
 \le(1+2d_0)\norm{T_{t^2}^G}_{2\to2}.
\]
It follows that
\[
 \ptwo(H)^2=\ptwo(G).
\]
For every $0\le p<\ptwo(G)$, apply the simple-graph exponential
bound to $H$ at $t=\sqrt p<\ptwo(H)$. Old-vertex connectivity agrees
and $\dist_H(x,y)=2\dist_G(x,y)$ for $x,y\in V(G)$, so restriction
gives \eqref{eq:exponential-connectivity} with constants depending
only on $G$ and $p$. Its validity for every $p<\ptwo(G)$ implies
$\ptwo(G)\le\pexp(G)$, completing the proof.
\end{proof}
